\section{Camera reference (R2)}\label{sec:camera}

As a reference we use the CRB of a camera with a Gaussian point-spread function of
$\sigma_{\mathrm{PSF}}=100$\,nm \cite{Balzarotti2017}. In the continuum limit without background
the bound is $\sigma_{\mathrm{PSF}}/\sqrt N$:
\pnum{camera_ideal_N400_nm}\src{camera_ideal_N400_nm}\,nm for $N=400$ and
\pnum{camera_sigma_nm}\src{camera_sigma_nm}\,nm for $N=1000$; this \emph{ideal} value is the
camera reference of Sec.~\ref{sec:iterative}. Pixelation costs little: a $9\times9$ window of
$100$\,nm pixels gives \pnum{camera_pixelated_9x9_N400_nm}\src{camera_pixelated_9x9_N400_nm}\,nm
at $N=400$ and \pnum{camera_pixelated_9x9_N1000_nm}\src{camera_pixelated_9x9_N1000_nm}\,nm at
$N=1000$.

The SBR convention matters for the camera. Balzarotti's $\sbr_c$ is the total signal over the
background \emph{per pixel} (Eq.~S61), so $\sbr_c=500$ on $9\times9$ pixels corresponds to a total
SBR of \pnum{camera_sbr_perpixel_to_total}\src{camera_sbr_perpixel_to_total}. With that
convention the pixelated camera reaches
\pnum{camera_perpixel_sbr500_N600_nm}\src{camera_perpixel_sbr500_N600_nm}\,nm at $N=600$ and
needs \pnum{camera_photons_for_5nm_perpixel_sbr500}\src{camera_photons_for_5nm_perpixel_sbr500}
photons for $5$\,nm; if $500$ were the total SBR it would reach
\pnum{camera_total_sbr500_N600_nm}\src{camera_total_sbr500_N600_nm}\,nm and need
\pnum{camera_photons_for_5nm_total_sbr500}\src{camera_photons_for_5nm_total_sbr500} photons. The
SBR values of MINFLUX (total/total) and of the camera (total/per pixel) are therefore not
directly comparable. For comparison, MINFLUX with $L=50$\,nm and $\sbr=10$ needs
\pnum{photons_for_5nm_L50_sbr10}\src{photons_for_5nm_L50_sbr10} photons for a $5$\,nm centre CRB
[Fig.~\ref{fig:scaling}(d)].
